\chapter{Conclusion}
\label{chap:conclusion}

This project implemented and compared two families of gradient inversion
attacks against Federated Learning's informal privacy claim --- optimization-based
(DLG, iDLG, Inverting Gradients) and closed-form (R-GAP) --- under four
controlled comparisons spanning architecture, scale, batch size, and a
DP-SGD defense (Chapter~\ref{chap:experiment}).

Across all four comparisons, a leaked gradient proved sufficient to
reconstruct private data under both attack families, confirming that the
privacy claim does not hold in general. The two families fail in
different ways: the closed-form attack is exact and fast but brittle,
degrading sharply under batch size and under noise, while the
optimization-based attacks are slower and less precise but degrade more
gracefully under the same conditions. Scaling to a deeper architecture is
only tractable for the optimization-based family, and at a cost in
fidelity. No single defense configuration protects against both families
equally; its effectiveness depends on which attack is actually in play.

Two extensions follow directly from this comparison. Pairing the
attack-side measurements of Section~\ref{sec:dpsgd} with a training-side
utility measurement would show whether a noise level sufficient to defeat
these attacks is one a real deployment could afford. Separately, R-GAP's
own rank-analysis framework (Section~\ref{sec:rgap}) could be used to
predict, for architectures not tested here (deeper convolutional networks,
BatchNorm, ReLU), whether a closed-form attack is even theoretically
possible before porting the attack itself.
